% !TEX root = main.tex

\section{\texorpdfstring{Distribution of the complexities of $\beta$-expansions}{}}

\label{sec:distribution of the complexities of beta expansions}

In Section \ref{sec:upper bound on the complexity of beta expansions}, we established an upper bound on the algorithmic complexity of the lexicographically largest elements of the equivalence classes generated by the equivalence relationship $\sim_\beta$ that naturally appears when $\beta$ is algebraic. In this section, we further study the algorithmic complexity of the remaining elements of those equivalence classes. When $\beta$ is Pisot, this allows us to characterize precisely the distribution of the algorithmic complexities of all the $\beta$-expansions of a given real number $s \in [0,1]$, in function of the algorithmic complexity of its binary expansions. We start by building a multivalued function that is computable relatively to $\beta$, that generates the set of all $\beta$-expansions. Later, we will use this function to link the algorithmic complexity of a given $\beta$-expansion to the algorithmic complexity of the sequence of coin tosses that is effected in line 6 of Algorithm \ref{alg:random algorithm beta expansion}. 

\subsection{\texorpdfstring{A computable function to generate all the $\beta$-expansions from a fixed $\beta$-expansion}{}}

Recall that for a fixed number $\beta \in (1,2)$, we have defined the equivalence relationship $\sim_\beta$ over $\{0,1\}^\ast$ by 
\begin{equation}
    x \sim_\beta y \iff n := \len x = |y| \ \text{and} \ \sum_{i=1}^n x_i \beta^{-i} = \sum_{i=1}^n y_i \beta^{-i}, \ \forall x,y \in \{0,1\}^\ast.
\end{equation}
Moreover, we have defined the equivalence class $[x]_\beta := \{y \in \{0,1\}^\ast : x \sim_\beta y\}$, for all $x \in \{0,1\}^\ast$, which, seen as a multivalued function $[\cdot]_\beta : \{0,1\}^\ast \rightrightarrows \{0,1\}^\ast$, is proven to be computable in the proof of Corollary \ref{cor:M is computable for algebraic beta}. 

We now construct a multivalued function $f_{\beta, 1 \to all} : \{0,1\}^\ast \rightrightarrows \{0,1\}^\ast$ that is computable relatively to $\beta$, such that
\begin{equation}
    \langle \Sigma_\beta(s,n) \rangle \in f_{\beta, 1 \to all}\left(x\right), \ \forall x \in \Sigma_\beta(s,n), \ s \in I_\beta, \ \text{and} \ n \in \N.
\end{equation}
To build $f_{\beta, 1 \to all}$, we first follow the same intuitions that lead to the construction of $f_{\beta \to 2}$ in Definition \ref{def:multivalued function converting from base beta to base 2} and $f_{2 \to \beta}$ in Definition \ref{def:multivalued function converting from base 2 to base beta}. Namely, the knowledge of $x \in \Sigma_\beta(s,n)$ allows to infer that
\begin{equation}
    s \in I(x) := \left[\sum_{i=1}^{n} x_i \beta^{-i},  \sum_{i=1}^{n} x_i \beta^{-i}+ \frac{\beta^{-n}}{\beta-1} \right],
\end{equation}
and further deduce that 
\begin{equation}
    \sum_{i=1}^n y_i \beta^{-i} \in J(x) := \left[\sum_{i=1}^{n} x_i \beta^{-i} - \frac{\beta^{-n}}{\beta-1},  \sum_{i=1}^{n} x_i \beta^{-i}+ \frac{\beta^{-n}}{\beta-1} \right],
\end{equation}
for all $y \in \Sigma_\beta(s,n)$. This leads to the definition of the following function.
\begin{definition}\label{def:definition of g beta to be all the beta sequences that are inside of an interval}
    For $\beta \in (1,2)$, define $g_\beta : \{0,1\}^\ast \rightrightarrows \{0,1\}^\ast$ as 
\begin{equation}
    g_\beta(x) := \left\{ y \in \{0,1\}^n : n = \len x,  \sum_{i=1}^n y_i \beta^{-i} \in_n J(x)\right\},
\end{equation}
for all $x \in \{0,1\}^\ast$.
\end{definition}
\noindent By the preceding discussion, we have
\begin{equation}\label{eq:Sigma beta s n is included in g beta x}
    \Sigma_\beta(s,n) \subseteq g_\beta(x), \ \forall x \in \Sigma_\beta(s,n).
\end{equation}
Now, define $\tilde g_\beta :  \{0,1\}^\ast \to \{0,1\}^\ast / \sim_\beta$ by
\begin{equation}
    \tilde g_\beta(x) := g_\beta(x)/ \sim_\beta, \ \forall x \in \{0,1\}^\ast,
\end{equation}
and
\begin{equation}
    \tilde \NN_{g_\beta}(x) := \# \tilde g_\beta(x), \ \forall x \in \{0,1\}^\ast.
\end{equation}
Note that the function defined by
\begin{equation}
    \begin{array}{rccc}
        \varphi_\beta : & \{0,1\}^n / \sim_\beta & \to & I_\beta\\
        & [x]_\beta & \mapsto & \sum_{i=1}^n x_i \beta^{-i}
    \end{array}
\end{equation}
is injective into $I_\beta$, since by definiton of $\sim_\beta$, $[x]_\beta \neq [y]_\beta \iff \sum_{i=1}^n x_i \beta^{-i} \neq \sum_{i=1}^n x_i \beta^{-i}$, for all $x,y \in \{0,1\}^n$. We define an ordering $\prec$ on $\{0,1\}^n / \sim_\beta$ by 
\begin{equation}
    u \prec v \iff \varphi_\beta(u) < \varphi_\beta(v), \ \forall u,v \in \{0,1\}^n / \sim_\beta.
\end{equation}
For $x \in \{0,1\}^\ast$, we enumerate the elements of $\tilde g_\beta(x)$ as follows. For $i \in \{1,\ldots, \tilde \NN_{g_\beta}(x)\}$, we let $\tilde g_\beta(x)|_i \in \tilde g_\beta(x)$, such that
\begin{equation}
    \tilde g_\beta(x)|_i < \tilde g_\beta(x)|_{i+1}, \ \forall i \in \{1,\ldots, \tilde \NN_{g_\beta}(x) - 1\}.
\end{equation}
We can then express $\Sigma_\beta(s,n)$ simply in function of the elements of $\tilde g_\beta(x)$. 
\begin{lemma}\label{lem:the set of beta expansion is expressed as a union of consecutive equivalence classes}
    Let $\beta \in (1,2)$, $s \in [0,1]$, $n \in \N$ and $x \in \Sigma_\beta(s,n)$. Then, there exists $i,j \in \{1, \ldots, \tilde \NN_{g_\beta}(x)\}$, $i \leq j$, such that 
    \begin{equation}
        \Sigma_\beta(s,n) = \bigcup_{\ell=i}^j\tilde g_\beta(x)|_\ell.
    \end{equation}
\end{lemma}
\begin{proof}
    Let $\beta \in (1,2)$, $s \in [0,1]$, $n \in \N$, and $x \in \Sigma_\beta(s,n)$. The proof is in two parts: we first prove that $x \in \Sigma_\beta(s,n) \Rightarrow [x]_\beta \subseteq \Sigma_\beta(s,n)$, and further that $[x]_\beta \prec [y]_\beta \prec [z]_\beta$ together with $x,z \in \Sigma_\beta(s,n)$ implies $y \in \Sigma_\beta(s,n)$. 
    \begin{enumerate}[label = (\alph*)]
        \item Let $x \in \Sigma_\beta(s,n)$. By Lemma \ref{lem:sigma beta s n contains the equivalence classes}, $[x]_\beta \in \Sigma_\beta(s,n)$. Hence, $\Sigma_\beta(s,n)$ is expressed as 
        \begin{equation}
            \Sigma_\beta(s,n) = \bigcup_{u \in X} u,
        \end{equation}
        for some subset $X \subseteq \{0,1\}^n/ \sim_\beta$. Since by (\ref{eq:Sigma beta s n is included in g beta x}), $\Sigma_\beta(s,n) \subseteq g_\beta(x)$, we get that $u \subseteq g_\beta(x) \Rightarrow u \in g_\beta(x)/ \sim_\beta \ = \tilde g_\beta(x)$, so there exists $\ell \in \{1,\ldots, \tilde\NN_{g_\beta} (x)\}$ such that $u = \tilde g_\beta(x)|_\ell$, for all $u \in X$. Therefore,
        \begin{equation}
            \Sigma_\beta(s,n) = \bigcup_{\ell \in L} \tilde g_\beta(x)|_\ell,
        \end{equation}
        for some $L \subseteq \{1, \ldots, \tilde \NN_{g_\beta}(x)\}$.
        \item We now end the proof by showing that $L$ is made of consecutive elements. Let $i,j,k \in \{1, \ldots, \tilde \NN_{g_\beta}(x)\}$, such that $i \leq j \leq k$ and suppose that $\tilde g_\beta(x)|_i, \tilde g_\beta(x)|_k \subseteq \Sigma_\beta(s,n)$. Then, for all $y \in u_j(x)$, all $u \in \tilde g_\beta(x)|_i$ and all $v \in \tilde g_\beta(x)|_k$, we have 
        \begin{equation}
            \sum_{m=1}^n u_m \beta^{-m} \leq \sum_{m=1}^n y_m \beta^{-m} \leq \sum_{m=1}^n v_m \beta^{-m}.
        \end{equation}
        Since $\tilde g_\beta(x)|_i, \tilde g_\beta(x)|_k \subseteq \Sigma_\beta(s,n)$, we further get
        \begin{equation}
            s - \frac{\beta^{-n}}{\beta-1} \leq \sum_{m=1}^n u_m \beta^{-m} \leq \sum_{m=1}^n y_m \beta^{-m} \leq \sum_{m=1}^n v_m \beta^{-m} \leq s.
        \end{equation}
        By (\ref{eq:1:construction of a multivalued function computable wrt beta that maps beta expansions to binary expansions of a real number}) and (\ref{eq:2:construction of a multivalued function computable wrt beta that maps beta expansions to binary expansions of a real number}), this implies that $y \in \Sigma_\beta(s,n)$, so $\tilde g_\beta(x)|_j \subseteq \Sigma_\beta(s,n)$. We have proven that $i,k \in L \Rightarrow j \in L$ for all $i \leq j \leq k$, i.e., $L$ is made of consecutive elements.
    \end{enumerate}
\end{proof}
\noindent For every $M \in \N$, we define $\CC(M)$ to be the set of all subset of consecutive elements in $\{1,\ldots,M\}$. In mathematical symbols, this is expressed as
\begin{equation}
    \CC(M) := \{ \{i, \ldots, j\} : i,j \in \{1,\ldots,M\}, i \leq j\}.
\end{equation}
We are ready to state the definition of the computable function $f_{\beta, 1 \to all}$. 
\begin{definition}\label{def:definition of f beta 1 all}
    Let $\beta \in (1,2)$. Define $\iota_\beta(x)$ to be the unique natural number such that 
    \begin{equation}
        x \in \tilde g_\beta(x)|_{\iota_\beta(x)}.
    \end{equation}
    Let $f_{\beta, 1 \to all} : \{0,1\}^\ast \rightrightarrows \{0,1\}^\ast$ be defined as
    \begin{equation}\label{eq:main:def:definition of f beta 1 all}
        f_{\beta, 1 \to all}(x) := \left\{ \left\langle \bigcup_{i \in A} \tilde g_{\beta}(x)|_i \right\rangle : A \in \CC\left(\tilde \NN_{g_\beta}(x)\right), \iota_\beta(x) \in A\right\},
    \end{equation}
    for all $x \in \{0,1\}^\ast$.
\end{definition}
\begin{lemma}
    Let $\beta \in (1,2)$, $s \in I_\beta$ and $n \in \N$. Then,
    \begin{equation}
        \langle \Sigma_\beta(s,n) \rangle \in f_{\beta, 1 \to all}(x), \ \forall x \in \Sigma_\beta(s,n).
    \end{equation}
\end{lemma}
\begin{proof}
    Let $\beta \in (1,2)$, $s \in I_\beta$, $n \in \N$ and $x \in \Sigma_\beta(s,n)$. Then, by Lemma \ref{lem:the set of beta expansion is expressed as a union of consecutive equivalence classes}, there exists $A \in \CC(\tilde \NN_{g_\beta}(x))$ such that 
    \begin{equation}
        \Sigma_\beta(s,n) = \bigcup_{i \in A} \tilde g_\beta(x)|_i.
    \end{equation}
    The Lemma follows immediately by Definition \ref{def:definition of f beta 1 all} of $f_{\beta, 1 \to all}$.
\end{proof}
% \begin{lemma}\label{lem:f beta 1 all is computable relatively to beta}
%     Let $\beta \in (1,2)$. Then, $f_{\beta, 1 \to all}$ is computable relatively to $\beta$.
% \end{lemma}
\begin{lemma}\label{lem:size of the multifunction that generates all beta expansions from one}
    Let $\beta \in (1,2)$ be an algebraic number and $n \in \N$. Then,
    \begin{equation}
        \#f_{\beta, 1 \to all}(x) \leq \left(\frac{1}{\beta-1} + 2\right)^2\frac{n^{2k_\beta}}{(L_\beta\Pi_\beta)^2}\left(L_\beta \Pi_\beta^+ \right)^{2n},
    \end{equation}
    for all $x \in \{0,1\}^n$.
\end{lemma}
\begin{proof}
    Let $\beta \in (1,2)$ be an algebraic number, $n \in \N$ and $x \in \{0,1\}^n$. By Lemma \ref{lem:generalization of the garsia separation lemma}, we get that 
    \begin{equation}
        |\varphi_\beta(u) - \varphi_\beta(v)| \geq \frac{L_\beta\Pi_\beta}{n^{k_\beta}\left(\beta L_\beta  \Pi_\beta^+ \right)^n} =: m_n,
    \end{equation}
    for all $u,v \in \{0,1\}^n / \sim_\beta$ with $u \neq v$. Then,
    \begin{align}
        \# \left(\varphi_\beta(\{0,1\}^n / \sim_\beta) \cap J(x)^{(n)}\right) &\leq \frac{|J(x)^{(n)}|}{m_n}\\
        & = \frac{2}{m_n}\left(\frac{\beta^{-n}}{\beta-1} + 2 \cdot 2^{-n}\right)\\
        & \leq \frac{2\beta^{-n}}{m_n}\left(\frac{1}{\beta-1} + 2\right) \cdot
    \end{align}
    By injectivity of $\varphi_\beta$,
    \begin{align}
        \# \varphi_\beta^{-1} \left(\varphi_\beta(\{0,1\}^n / \sim_\beta) \cap J(x)^{(n)}\right) &= \# \left(\varphi_\beta(\{0,1\}^n / \sim_\beta) \cap J(x)^{(n)}\right)\\ 
        &\leq \frac{2\beta^{-n}}{m_n}\left(\frac{1}{\beta-1} + 2\right) \cdot
    \end{align}
    Hence, since by Definition \ref{def:definition of g beta to be all the beta sequences that are inside of an interval}, $\tilde g_\beta(x) \subseteq \varphi_\beta^{-1} \left(\varphi_\beta(\{0,1\}^n / \sim_\beta) \cap J(x)^{(n)}\right)$ for $x \in \{0,1\}^n$, we get that 
    \begin{equation}
        \tilde \NN_{g_\beta}(x) \leq  \frac{2\beta^{-n}}{m_n}\left(\frac{1}{\beta-1} + 2\right) = \left(\frac{1}{\beta-1} + 2\right)\frac{2n^{k_\beta}}{L_\beta \Pi_\beta} \left(L_\beta \Pi_\beta^+ \right)^{n}.
    \end{equation}
    
    \noindent Now, there remains to establish a bound on $\#f_{\beta, 1 \to all}(x) = \#\{ A : A \in \CC(\tilde \NN_{g_\beta}(x)), \iota_\beta(x) \in A\}$, for $x \in \{0,1\}^n$. Let $M \in \N$, $i \in \{1,\ldots, i\}$, and define the set $\DD(M,i) := \{ A : A \in \CC(M), i \in A\}$. We will establish a general bound on $\#\DD(M,i)$, which as a corollary will give a bound on $\#\DD(\tilde \NN_{g_\beta}(x),\iota_\beta(x))$. By symmetry, we can restrict ourselves to $i \leq M/2$ without loss of generality. Define the sets $\DD_k(M,i) = \{ A \in \DD(M,i) : \#A = k\}$, for $k \in \{1, \ldots, M\}$. Note that 
    \begin{equation}
        \#\DD_k(M,i) = \min\{k, i, M-k\}.
    \end{equation}
    Then,
    \begin{align}
        \# \DD(M,i) &= \sum_{k = 1}^M \#\DD_k(M,i) = \sum_{k = 1}^M \min\{k, i, M-k\}\\
        &= \sum_{k=1}^{i-1} k + \sum_{k=i}^{M-i} i + \sum_{k=M - i +1}^{M} (M -k)\\
        & = 2 \sum_{k=1}^{i-1} + (M - 2i + 1)i = i(i-1) + (M - 2i + 1)i\\
        &= Mi - i^2.
    \end{align}
    Note that $Mi - i^2$ is maximal for $i = M/2$. Therefore,
    \begin{equation}
        \# \DD(M,i) \leq \frac14 M^2.
    \end{equation}
    We conclude that 
    \begin{align}
        \#f_{\beta, 1 \to all}(x) &= \#\DD(\tilde \NN_{g_\beta}(x),\iota_\beta(x)) \leq \frac14 \tilde \NN^2_{g_\beta}(x)\\
        & \leq \left(\frac{1}{\beta-1} + 2\right)^2\frac{n^{2k_\beta}}{((\beta-1)L_\beta\Pi_\beta)^2}\left(L_\beta \Pi_\beta^+ \right)^{2n},
    \end{align}
    for all $x \in \{0,1\}^n$.
\end{proof}

\subsection{\texorpdfstring{An algorithm to extract the coin tosses from a $\beta$-expansion}{}}

In this part, we show that every $\beta$-expansion of a given number $s \in I_\beta$ can be associated with a sequence of coin tosses generated by Algorithm \ref{alg:random algorithm beta expansion}. We will then study the algorithmic complexity of $\beta$-expansions through the study of this sequence of coin tosses. Following \cite[Section 1.2]{dajani2003random}, we first establish the equations that govern Algorithm \ref{alg:random algorithm beta expansion}, given fixed input $\beta \in (1,2)$, $s \in [0,1]$, and $n \in \N$. First remark that Algorithm \ref{alg:random algorithm beta expansion} generates two sequences $(b_i)_{i \in \{1,\ldots,n\}}$ and $(r_i)_{i \in \{0,\ldots,n\}}$, that are recursively defined according to the following pattern. By lines 1 and 9 in Algorithm \ref{alg:random algorithm beta expansion}, we get that 
\begin{equation}
    r_0 = s, \ \ r_i = \beta r_{i-1} - b_i, \ \forall i \in \{1, \ldots, n\}.
\end{equation}
Further, lines 4-6 deliver
\begin{equation}
    b_i = \begin{cases}
        0 \ \text{if} \ r_{i-1} \in [0, \beta^{-1}) =: E_\beta^{(0)}\\
        \text{Randomly} \ 0 \text{ or } 1 \ \text{if} \  r_{i-1}  \in [\beta^{-1}, ((\beta-1)\beta)^{-1}] =: S_\beta\\
        1 \ \text{if} \ r_{i-1} \in (((\beta-1)\beta)^{-1}, (\beta-1)^{-1}] =: E_\beta^{(1)}
    \end{cases} =: T_\beta(r_{i-1}).
\end{equation}
Note that the recursion generating $(r_i)_{i \in \{1, \ldots , n\}}$ can be simply expressed by 
\begin{equation}\label{eq:definition of K beta, simple recursive relation for r i}
    r_i = \beta r_{i-1} - T_\beta(r_{i-1}) =: K_\beta(r_{i-1}), \ \forall i \in \{1, \ldots, n\}.
\end{equation}
The output $\AM_\beta (s,n)$ of the algorithm is given in line 8 by 
\begin{equation}
    \AM_\beta (s,n) = \coprod_{i=1}^n b_i = \coprod_{i=1}^n T_\beta(r_{i-1}) =  \coprod_{i=0}^{n-1} T_\beta(r_i).
\end{equation}
In order to formalize better the problem, we can define $T_\beta(r)$ formally as a random variable on $\{0,1\}$ for $r \in I_\beta$, i.e., as a function $T_\beta(r) : \{0,1\} \to \{0,1\}$ satisfying
\begin{equation}
    T_\beta(r)(x) := \begin{cases}
        0, \ \text{if} \ s \in  E_\beta^{(0)},\\
        x, \ \text{if} \ s \in  S_\beta,\\
        1, \ \text{if} \ s \in  E_\beta^{(1)},
    \end{cases}
    \ \forall x \in \{0,1\}, \ r \in I_\beta.
\end{equation}
Since the generation of the sequence $(r_i)_{i \in \N}$ might rely on an arbitrarily many realisations of the above random variable, we define an extension $\tilde T_\beta$ of $T_\beta$, that will be useful to consider Algorithm \ref{alg:random algorithm beta expansion} as a whole to be a random variable. Namely, we define $\tilde T_\beta(r) : \{0,1\}^\N \to \{0,1\}$ by 
\begin{equation}\label{eq:definition of the generalized random bit generator}
    \tilde T_\beta(r)(\xf) := \begin{cases}
        0, \ \text{if} \ s \in  E_\beta^{(0)},\\
        \xf_1, \ \text{if} \ s \in  S_\beta,\\
        1, \ \text{if} \ s \in  E_\beta^{(1)},
    \end{cases}
    \ \forall \xf \in \{0,1\}^\N, \ r \in I_\beta.
\end{equation}
Thanks to the formulation above, we can see Algorithm \ref{alg:random algorithm beta expansion} as first generating an infinite sequence $\xf \in \{0,1\}^\N$ out of infinitely many Bernoulli experiments, and then reading successively the outcomes of this experiment, i.e., at each time step, the algorithm reads $\xf_1$, and then discards it, such that the sequence $\xf$ is transformed into the sequence $\xf' := \xf_2\xf_3\ldots$. This is usually denoted by $\xf' = \sigma(\xf)$, where $\sigma : \{0,1\}^\N \to \{0,1\}^\N$ is called the left-shift. Then, in the next time step, the algorithm will read $\xf'_1 = \xf_2$ to generate the next bit of the $\beta$-expansion of $s$.  Note that when $r_i \notin S_\beta$, then $\tilde T_\beta(r_i)(\xf)$ does not depend on $\xf \in \{0,1\}^\N$, as it is generated deterministically. Hence, the Algorithm \ref{alg:random algorithm beta expansion} does not need to discard $\xf_1$, and wait until the next time step $i$ for which $r_i \in S_\beta$ to discard $\xf_1$. Accordingly, we define $\sigma_r : \{0,1\}^\N \to \{0,1\}^\N$ by 
\begin{equation}
    \sigma_r(\xf) := \begin{cases}
        \xf \ \text{if} \ r \notin S_\beta,\\
        \sigma(\xf), \ \text{if} \ r \in S_\beta,
    \end{cases} \ \forall r \in I_\beta, \ \xf \in \{0,1\}^\N.
\end{equation}
We are now ready to state formally the definition of the Algorithm \ref{alg:random algorithm beta expansion} as a random variable. First, we model $K_\beta(r)$ defined in (\ref{eq:definition of K beta, simple recursive relation for r i}) as a random variable on $\{0,1\}^\N$ for $r \in I_\beta$, i.e., as a function $K_\beta(r) : \{0,1\}^\N \to I_\beta \times \{0,1\}^\N$,  defined by 
\begin{equation}
    K_\beta(r)(\xf) := (\beta r - \tilde T_\beta(r)(\xf_1), \sigma_r(\xf)),
    \ \forall \xf \in \{0,1\}^\N.
\end{equation}
Note that thanks to the formalism introduced, we can also consider the sequences $(b_i)_{i \in \{1, \ldots, n\}}$ and $(r_i)_{i \in \{0, \ldots, n\}}$ as random variables, i.e., each elements is a function from $\{0,1\}^\N$ to $\{0,1\}$ and $I_\beta$, respectively. We denote their respective elements $b_i(\xf)$ and $r_i(\xf)$, for $\xf \in \{0,1\}^\N$. In order to simplify the notations, we might consider $\tilde T_\beta$ as a function from $I_\beta \times \{0,1\}^\N$ into $\{0,1\}$ and $K_{\beta}$ as a function from $I_\beta \times \{0,1\}^\N$ into itself, by making the confusion between $\tilde T_\beta(s)(\xf)$ and $\tilde T_\beta(s,\xf)$, and the confusion between $K_{\beta}(s)(\xf)$ and $K_{\beta}(s,\xf)$. Note that this way, we get a simple equation for $r_i(\xf)$ in function of $\xf$:
\begin{equation}
    r_i(\xf) = \pi_1 K^i_\beta(s,\xf), \ \forall \xf \in \{0,1\}^\N,
\end{equation}
where $\pi_1$ is the projection onto the first coordinate. This simpler formalism allows to define the output of the Algorithm \ref{alg:random algorithm beta expansion} as a random variable as well, i.e., as a function $\AM_\beta(s,n) : \{0,1\}^\N \to \{0,1\}^n$, by 
\begin{equation}\label{eq:definition of the random beta expansion algorithm as a random variable}
    \AM_\beta(s,n)(\xf) = \coprod_{i=0}^{n-1} \tilde T_\beta\left(K_{\beta}^i(s,\xf)\right), \ \forall \xf \in \{0,1\}^\N.
\end{equation}

We repeat that Algorithm \ref{alg:random algorithm beta expansion} can generate different outputs of the same input $s \in [0,1]$. \cite[Theorem 1]{dajani2007invariant} establishes that infact, every $\beta$-expansion of $s$ can actually be generated by this Algorithm. This is formally expressed by 
\begin{equation}
    \AM_\beta(s,n) \left( \{0,1\}^\N\right) = \Sigma_\beta(s,n).
\end{equation}
We let also $\NN_\beta(s,n) :=  \# \left(\Sigma_\beta(s,n) \right)$.
We now leverage the formalism introduced above to express every $\beta$-expansion as a decomposition into a binary expansion and a sequence of coin tosses. We start by showing that for every $\beta$-expansion of $s$, we can associate a sequence of coin tosses that generated it. First, define
\begin{equation}
    \SM_\beta(s,\xf) := \{i \in \N_0 : \pi_1 K_\beta^i(s,\xf) \in S_\beta\},
\end{equation}
\begin{equation}
    \SM_\beta(s,n,\xf) := \SM_\beta(s,\xf) \cap \{0,\ldots,n-1\},
\end{equation} 
and
\begin{equation}
    h_\beta(s,n,\xf) := \# \SM_\beta(s,n,\xf) \leq n,
\end{equation}
for all $\xf \in \{0,1\}^\N$. We first note a simple technical Lemma.
\begin{lemma}\label{lem:recursive relation for determining h beta}
    Let $\beta \in (1,2)$, $s \in I_\beta$, $n \in \N$, and $\xf \in \{0,1\}^\N$. Then,
    \begin{equation}\label{eq:main1:lem:recursive relation for determining h beta}
        h_\beta(s,n+1,\xf) = h_\beta(s,n, \xf) + 1_{n \in \SM_\beta(s,\xf)},
    \end{equation}
    and therefore,
    \begin{equation}\label{eq:main2:lem:recursive relation for determining h beta}
        \bigcup_{i \in \SM_\beta(s,n,\xf)} \{h_\beta(s,i+1,\xf)\} = \{1, \ldots, h_\beta(s,n,\xf)\}.
    \end{equation}
\end{lemma}
\begin{proof}
    Let $\beta \in (1,2)$, $s \in I_\beta$, $n \in \N$, and $\xf \in \{0,1\}^\N$. Suppose that $n \notin S_\beta(s,\xf)$. Then,
    \begin{equation}
        \SM_\beta(s,n+1,\xf) = \SM_\beta(s,\xf) \cap \{0,\ldots, n\} = \SM_\beta(s,\xf) \cap \{0,\ldots, n-1\} = \SM_\beta(s,n,\xf),
    \end{equation}
    therefore $h_\beta(s,n+1,\xf) = h_\beta(s,n, \xf)$. Suppose now that $n \in S_\beta(s,\xf)$. Then,
    \begin{align}
        \SM_\beta(s,n+1,\xf) &= \SM_\beta(s,\xf) \cap \{0,\ldots, n\}\\
        & = \left(\SM_\beta(s,\xf) \cap \{0,\ldots, n-1\}\right) \cup \{n\}\\
        & = \SM_\beta(s,n,\xf) \cup \{n\}\,
    \end{align}
    therefore $h_\beta(s,n+1,\xf) = h_\beta(s,n, \xf) + 1$. This yields (\ref{eq:main1:lem:recursive relation for determining h beta}). Let $i_1 < i_2 < \ldots < i_{h_\beta(s,n, \xf)}$ be the elements of $\SM_\beta(s,n,\xf)$. For all $j \leq i_1$, $h_\beta(s, j, \xf) = h_\beta(s, 0, \xf) = 0$, so $h_\beta(s, i_1+1, \xf) = h_\beta(s, i_1, \xf) + 1 = 1$. Moreover, for all $k \in \{2, \ldots, h_\beta(s,n, \xf)\}$,
    \begin{equation}
        h_\beta(s, i_k+1, \xf) = h_\beta(s, i_k, \xf) + 1 = h_\beta(s, i_{k-1}+1, \xf) + 1,
    \end{equation}
    which by induction on $k$, yields
    \begin{equation}
        h_\beta(s, i_k+1, \xf) = k.
    \end{equation}
    Then,
    \begin{align}
        \bigcup_{i \in \SM_\beta(s,n,\xf)} \{h_\beta(s,i+1,\xf)\} &= \bigcup_{k = 1}^{h_\beta(s,n,\xf)} \{h_\beta(s,i_k+1,\xf)\}\\
        & = \bigcup_{k = 1}^{h_\beta(s,n,\xf)} \{k\} \\
        &= \{1, \ldots, h_\beta(s,n,\xf)\}.
    \end{align}
\end{proof}
\noindent $h_\beta(s,n-1,\xf)$ is exactly the number of bits of $\xf$ that Algorithm \ref{alg:random algorithm beta expansion} had to read before delivering its output. Then, every bit of $\xf$ that is localed after the rank $h_\beta(s,n,\xf)$ is not read by the algorithm to produce its output. This is formally expressed in the following Lemma and the subsequent Corollary.
\begin{lemma}\label{lem:the random beta expansion algorithm reads only the first bits of its input}
    Let $\beta \in (1,2)$, $s \in I_\beta$, $n \in \N$. Then,
    \begin{equation}
        \begin{cases}
            \SM_\beta(s,n,\xf) = \SM_\beta(s,n,\xf_{1:h_\beta(s,n,\xf)}\yf),\\
            \displaystyle \pi_1 K^n_\beta(s,\xf) = \pi_1 K^n_\beta\left(s,\xf_{1:h_\beta(s,n,\xf)}\yf\right),\\
            \displaystyle \pi_2 K^n(s,\xf) = \sigma^{h_\beta(s,n,\xf)}(\xf),
        \end{cases}
    \end{equation}
    for all $\xf, \yf \in \{0,1\}^\N$, where $\pi_1$ (resp. $\pi_2$) is the projection on the first coordinate (resp. second coordinate).
\end{lemma}
\begin{proof}
    Let $\beta \in (1,2)$, $s \in I_\beta$. We prove the result by induction. Obviously,
    \begin{equation}
        \SM_\beta(s,0,\xf) = \varnothing = \SM_\beta(s,0,\xf_{1:h_\beta(s,0,\xf)}\yf),
    \end{equation}
    \begin{equation}
        \pi_1 K^0_\beta(s,\xf) = s =  \pi_1 K^0_\beta(s,\xf_{1:h_\beta(s,0,\xf)}\yf),
    \end{equation}
    and
    \begin{equation}
        \pi_2 K^i(s,\xf) = \xf = \sigma^{0}(\xf) = \sigma^{h_\beta(s,0,\xf)}(\xf),
    \end{equation}
    for all $\xf,\yf \in \{0,1\}^\N$.
    Now, suppose that there exists $n \in \N$ such that
    \begin{equation}\tag{$H_n$}
        \begin{cases}
            \SM_\beta(s,n,\xf) = \SM_\beta(s,n,\xf_{1:h_\beta(s,n,\xf)}\yf) =: \SM_n,\\
            \displaystyle \pi_1 K^n_\beta(s,\xf) = \pi_1 K^n_\beta\left(s,\xf_{1:h_\beta(s,n,\xf)}\yf\right) =: r_n,\\
            \displaystyle \pi_2 K^n(s,\xf) = \sigma^{h_\beta(s,n,\xf)}(\xf) = \sigma^{h_n}(\xf),
        \end{cases}
    \end{equation}
    where $h_n := \# \SM_n$, for all $\xf, \yf \in \{0,1\}^\N$. We will show that $(H_{n+1})$ holds for all $\xf, \yf \in \{0,1\}^\N$. Fix $\xf, \yf \in \{0,1\}^\N$. First, note that 
    \begin{align}
        K^{n+1}_\beta(s,\xf) &= K_\beta \left( K^n_\beta(s,\xf)\right)\\
        &= K_\beta\left( \pi_1 K^n_\beta(s,\xf), \pi_2 K^n_\beta(s,\xf)\right)\\
        &\overset{(a)}{=} K_\beta\left( r_n, \ \sigma^{h_n}(\xf)\right),
    \end{align}
    and similarly
    \begin{align}
        K^{n+1}_\beta(s,\xf_{1:h_\beta(s,n+1,\xf)}\yf) &= K_\beta \left( K^n_\beta(s,\xf_{1:h_n}\yf)\right)\\
        &= K_\beta\left( \pi_1 K^n_\beta(s,\xf_{1:h_n}\yf), \pi_2 K^n_\beta(s,\xf_{1:h_n}\yf)\right)\\
        &\overset{(b)}{=} K_\beta\left(r_n, \ \sigma^{h_\beta(s,n,\xf_{1:h_n}\yf)}(\xf_{1:h_n}\yf)\right),\\
        &\overset{(c)}{=} K_\beta\left(r_n, \ \sigma^{h_n}(\xf_{1:h_n}\yf)\right),
    \end{align}
    where (a), (b) and (c) follow from $(H_n)$.
    \begin{enumerate}[label = (\alph*)]
        \item Suppose that $r_n \in E_\beta^{(0)} \cup E_\beta^{(1)}$, and define $b \in \{0,1\}$ such that $r_n \in E_\beta^{(b)}$. Then, 
            \begin{enumerate}[label=(\roman*)]
                \item Since $r_n \notin S_\beta$, then $n \notin \SM_\beta(s,\xf)$, so
                \begin{align}
                    \SM_\beta(s,n+1,\xf) &= \SM_\beta(s,\xf) \cap \{0, \ldots, n\}\\
                    & =  \SM_\beta(s,\xf) \cap \{0, \ldots, n-1\}\\
                    & = \SM_\beta(s,n,\xf) = \SM_n,
                \end{align}
                and similarly, $n \notin \SM_\beta(s,\xf_{1:h_n}\yf)$, so
                \begin{align}
                    \SM_\beta(s,n+1,\xf_{1:h_n}\yf) &= \SM_\beta(s,\xf_{1:h_n}\yf) \cap \{0, \ldots, n\}\\
                    & =  \SM_\beta(s,\xf_{1:h_n}\yf) \cap \{0, \ldots, n-1\}\\
                    & = \SM_\beta(s,n,\xf_{1:h_n}\yf) = \SM_n,
                \end{align}
                which yields
                \begin{equation}
                    \SM_\beta(s,n+1,\xf) = \SM_\beta(s,n+1,\xf_{1:h_\beta(s,n,\xf)}\yf) =: \SM_{n+1},
                \end{equation}
                and note that $\SM_{n+1} = \SM_n$, and hence $h_{n+1} = h_n$.
                \item Note that
                \begin{equation}
                    \pi_1 K_\beta^{n+1}(s,\xf) = \pi_1 K_\beta(r_n, \sigma^{h_n}(\xf)) = \beta r_n - b,
                \end{equation}
                and 
                \begin{align}
                    \pi_1 K_\beta^{n+1}(s,\xf_{1:h_n}\yf) = \pi_1 K_\beta(r_n, \sigma^{h_n}(\xf_{1:h_n}\yf)) = \beta r_n - b,
                \end{align}
                so 
                \begin{equation}
                    r_{n+1} := \pi_1 K_\beta^{n+1}(s,\xf) = \pi_1 K_\beta^{n+1}(s,\xf_{1:h_n}\yf),
                \end{equation}
                and note that $r_{n+1} = \beta r_n - b$.
                \item Since $r_n \notin S_\beta$, then 
                \begin{equation}
                    \pi_2 K_\beta^{n+1}(s,\xf) = \pi_2 K_\beta(r_n, \sigma^{h_n}(\xf)) = \sigma^{h_n}(\xf) = \sigma^{h_{n+1}}(\xf).
                \end{equation}
            \end{enumerate}
                \item Suppose that $r_n \in S_\beta$. Then,
                \begin{enumerate}[label=(\roman*)]
                    \item Since $r_n \in S_\beta$, then $n \in \SM_\beta(s,\xf)$, so
                    \begin{align}
                        \SM_\beta(s,n+1,\xf) &= \SM_\beta(s,\xf) \cap \{0, \ldots, n\}\\
                        & =  \left(\SM_\beta(s,\xf) \cap \{0, \ldots, n-1\}\right) \cup \{n\}\\
                        & = \SM_\beta(s,n,\xf) \cup \{n\} = \SM_n \cup \{n\},
                    \end{align}
                    and similarly, $n \notin \SM_\beta(s,\xf_{1:h_n}\yf)$, so
                    \begin{align}
                        \SM_\beta(s,n+1,\xf_{1:h_n}\yf) &= \SM_\beta(s,\xf_{1:h_n}\yf) \cap \{0, \ldots, n\}\\
                        & =  \left(\SM_\beta(s,\xf_{1:h_n}\yf) \cap \{0, \ldots, n-1\} \right) \cup \{n\}\\
                        & = \SM_\beta(s,n,\xf_{1:h_n}\yf) \cup \{n\} = \SM_n \cup \{n\},
                    \end{align}
                    which yields
                    \begin{equation}
                        \SM_\beta(s,n+1,\xf) = \SM_\beta(s,n+1,\xf_{1:h_n}\yf) =: \SM_{n+1},
                    \end{equation}
                    and note that $\SM_{n+1} = \SM_n \cup \{n\}$, and hence $h_{n+1} = h_n + 1$.
                    \item Note that
                    \begin{align}
                        \pi_1 K_\beta^{n+1}(s,\xf) &= \pi_1 K_\beta(r_n, \sigma^{h_n}(\xf))\\
                        & = \beta r_n - \left[\sigma^{h_n}(\xf)\right]_1\\
                        & = \beta r_n - \xf_{h_n+1},
                    \end{align}
                    and 
                    \begin{align}
                        \pi_1 K_\beta^{n+1}(s,\xf_{1:h_{n+1}}\yf) &= \pi_1 K_\beta(r_n, \sigma^{h_n}(\xf_{1:h_{n+1}}\yf))\\
                        & = \beta r_n - \left[\sigma^{h_n}(\xf_{1:h_{n+1}}\yf)\right]_1\\
                        & = \beta r_n - \left[\sigma^{h_n}(\xf_{1:h_n+1}\yf)\right]_1\\
                        & = \beta r_n - \xf_{h_n+1},
                    \end{align}
                    so 
                    \begin{equation}
                        r_{n+1} := \pi_1 K_\beta^{n+1}(s,\xf) = \pi_1 K_\beta^{n+1}(s,\xf_{1:h_{n+1}}\yf),
                    \end{equation}
                    and note that $r_{n+1} = \beta r_n - \xf_{h_n+1}$.
                    \item Since $r_n \in S_\beta$, then 
                    \begin{align}
                        \pi_2 K_\beta^{n+1}(s,\xf) &= \pi_2 K_\beta(r_n, \sigma^{h_n}(\xf)) = \sigma\left(\sigma^{h_{n}}(\xf)\right)\\
                        & = \sigma^{h_{n}+1}(\xf) = \sigma^{h_{n+1}}(\xf).
                    \end{align}
                \end{enumerate}
    \end{enumerate}
    We hence have proven that $(H_n) \Rightarrow (H_{n+1})$. We have completed the induction, thereby finishing the proof.
\end{proof}

\begin{corollary}Let $\beta \in (1,2)$, $s \in I_\beta$, $n \in \N$. Then,
    \begin{equation}\label{eq:dajani algorithm only depends on the bits that were read by the algorithm}
        \AM_\beta(s,n)(\xf_{1:h_\beta(s,n,\xf)}\{0,1\}^\N) = \{\AM_\beta(s,n)(\xf)\}
    \end{equation}
    for all $\xf \in \{0,1\}^\N$.
\end{corollary}
\begin{proof}
    Let $\beta \in (1,2)$, $s \in I_\beta$, $n \in \N$, $\xf\in \{0,1\}^\N$ and $\zf \in \xf_{1:h_\beta(s,i,\xf)} \{0,1\}^\N$, Note that there exists $\yf \in \{0,1\}^\N$ such that $\zf := \xf_{1:h_\beta(s,i,\xf)}\yf$. Fix $i \in \{0, \ldots, n-1\}$. By Lemma \ref{lem:the random beta expansion algorithm reads only the first bits of its input}, we have 
    \begin{equation}
        \begin{cases}
            h_\beta(s,i,\xf) = h_\beta(s,i,\zf) =: h_i,\\
            \displaystyle \pi_1 K^n_\beta(s,\xf) = \pi_1 K^n_\beta\left(s,\zf\right) =: r_i,\\
            \displaystyle \pi_2 K^i(s,\xf) = \sigma^{h_\beta(s,i,\xf)}(\xf) = \sigma^{h_i}(\xf),\\
            \displaystyle \pi_2 K^i(s,\zf) = \sigma^{h_\beta(s,i,\zf)}(\zf) = \sigma^{h_i}(\zf),
        \end{cases}
    \end{equation}
    Then,
    \allowdisplaybreaks
    \begin{align}
        \tilde T_\beta \left( K_\beta^i(s,\xf) \right) &= \tilde T_\beta \left(  \pi_1 K_\beta^i(s,\xf), \pi_2 K_\beta^i(s,\xf)\right) = \tilde T_\beta \left(  r_i, \sigma^{h_i}(\xf)\right)\\
        &= \nonumber \begin{cases}
            0 \text{ if } r_i \in E_\beta^{(0)}\\
            \left[\sigma^{h_i}(\xf)\right]_1 \text{ if } r_i \in S_\beta\\
            1 \text{ if } r_i \in E_\beta^{(1)}
        \end{cases} = \begin{cases}
            0 \text{ if } r_i \in E_\beta^{(0)}\\
            \xf_{h_i + 1} \text{ if } r_i \in S_\beta\\
            1 \text{ if } r_i \in E_\beta^{(1)}
        \end{cases} \overset{(a)}{=} \begin{cases}
            0 \text{ if } r_i \in E_\beta^{(0)}\\
            \xf_{h_{i+1}}\text{ if } r_i \in S_\beta\\
            1 \text{ if } r_i \in E_\beta^{(1)}
        \end{cases}
    \end{align}
    where (a) follows from (\ref{eq:main1:lem:recursive relation for determining h beta}). Note that $h_{i+1} = h_\beta(s,i+1,\xf) \leq h_\beta(s,n,\xf)$, so $\xf_{h_{i+1}} = \zf_{h_{i+1}}$. Therefore,
    \begin{align}
        \tilde T_\beta \left( K_\beta^i(s,\xf) \right) = \begin{cases}
        0 \text{ if } r_i \in E_\beta^{(0)}\\
        \zf_{h_{i+1}}\text{ if } r_i \in S_\beta\\
        1 \text{ if } r_i \in E_\beta^{(1)}
    \end{cases} = \tilde T_\beta \left(  r_i, \sigma^{h_i}(\zf)\right) = \tilde T_\beta \left( K_\beta^i(s,\zf) \right).
    \end{align}
    We conclude the proof by noting that 
    \begin{equation}
        \AM_\beta(s,n)(\xf) = \coprod_{i=0}^{n-1} \tilde T_\beta \left( K_\beta^i(s,\xf) \right) = \coprod_{i=0}^{n-1} \tilde T_\beta \left( K_\beta^i(s,\zf) \right) = \AM_\beta(s,n)(\zf),
    \end{equation}
    for all $\zf \in \xf_{1:h_\beta(s,n,\xf)}\{0,1\}^\N$.
\end{proof}
% From these observations, we establish the following lemma.
% \begin{lemma}
%     Let $\beta \in (1,2)$, $s \in I_\beta$ and $u \in \{0,1\}^\ast$. Then, for every $\xf,\yf \in u \{0,1\}^\N$,
%     \begin{equation}
%         K_\beta^i(s,\xf) = K_\beta^i(s,\yf), 
%     \end{equation}
%     for all $i \leq \min \left\{ \max \SM_\beta(s,|u|,\xf), \max \SM_\beta(s,|u|,\yf) \right\}$.
% \end{lemma}
% NE SERT A RIEN
% \begin{corollary}
%     Let $\beta \in (1,2)$, $s \in I_\beta$ and $u \in \{0,1\}^\ast$. Then, for every $\xf,\yf \in u \{0,1\}^\N$, 
%     \begin{equation}
%         \SM_\beta(s,|u|+1,\xf) = \SM_\beta(s,|u|+1,\yf),
%     \end{equation}
%     and therefore 
%     \begin{equation}
%         h_\beta(s,|u|+1,\xf) = h_\beta(s,|u|+1,\yf).
%     \end{equation}
% \end{corollary}
% \begin{proof}
%     Let $\beta \in (1,2)$, $s \in I_\beta$, $u \in \{0,1\}^\ast$ and $\xf,\yf \in u \{0,1\}^\N$. Then,
%     \begin{align}
%         \SM_\beta(s,|u|+1,\xf) &= \left\{ i \in \{0,\ldots, |u|\} : K_\beta^i(s,\xf) \in S_\beta\right\}\\
%         &\overset{(a)}{=} \left\{ i \in \{0,\ldots, |u|\} : K_\beta^i\left(s,\xf_{1:h_\beta(s,i,\xf)}\zf^{(i)}\right) \in S_\beta\right\},
%     \end{align}
%     for all $\zf^{(i)} \in \{0,1\}^\N$ and all $i \in \{0,\ldots, |u|\}$, where (a) follows from previous LEMMA. In particular, since $h_\beta(s,i,\xf) \leq i \leq |u|$, we can choose $\zf^{(i)}$ so that $\xf_{1:h_\beta(s,i,\xf)}\zf^{(i)} = \yf$, for all $i \in \{0,\ldots, |u|\}$. Together with the ABOVE EQUATION, this yields
%     \begin{equation}
%         \SM_\beta(s,|u|+1,\xf) = \left\{ i \in \{1,\ldots, |u|\} : K_\beta^i(s,\yf) \in S_\beta\right\} = \SM_\beta(s,|u|+1,\yf).
%     \end{equation}
% \end{proof}
We now show that for every prefix $x$ of some $\beta$-expansion of $s$, is generated by Algorithm \ref{alg:random algorithm beta expansion} through a unique sequence of coin tosses denoted $w_\beta(s,x)$.
\begin{lemma}\label{lem:the set of preimages by the dajani algorithm is a cylinder}
    Let $\beta \in (1,2)$, $s \in I_\beta$ and $n \in \N$. Then, for every $x \in \Sigma_\beta(s,n)$, there exists a unique $w_\beta(s, x) \in \{0,1\}^{\leq n}$, such that 
    \begin{equation}
        \AM_\beta(s,n)^{-1}(x) = w_\beta(s, x) \{0,1\}^\N.
    \end{equation}
    Moreover, 
    \begin{equation}
        \xf_{1:h_\beta(s,n,\xf)} = w_\beta(s,x), \ \forall \xf \in \AM_\beta(s,n)^{-1}(x).
    \end{equation}
\end{lemma}
\begin{proof}
    Let $\beta \in (1,2)$, $s \in I_\beta$, $n \in \N$ and $x \in \Sigma_\beta(s,n)$. We first prove that $\AM_\beta(s,n)^{-1}(x) = A_x \{0,1\}^\N$, where $A_x \subseteq \{0,1\}^{\leq n}$, and then that $A_x$ consists of a unique element. 
    \begin{enumerate}[label = (\alph*)]
        \item Let $\xf \in \AM_\beta(s,n)^{-1}(x)$. Then, by (\ref{eq:dajani algorithm only depends on the bits that were read by the algorithm}), every $\yf \in \xf_{1:h_\beta(s,n,\xf)} \{0,1\}^\N$ satisfies $\yf \in \AM_\beta(s,n)^{-1}(x)$. It follows that
        \begin{equation}\label{eq:1:proof:lem:the set of preimages by the dajani algorithm is a cylinder}
            \AM_\beta(s,n)^{-1}(x) = \bigcup_{\xf \in \AM_\beta(s,n)^{-1}(x)} \xf_{1:h_\beta(s,n,\xf)} \{0,1\}^\N.
        \end{equation}
        By defining $A_x := \bigcup_{\xf \in \AM_\beta(s,n)^{-1}(x)} \{\xf_{1:h_\beta(s,n,\xf)}\}$, we then get 
        \begin{equation}
            \AM_\beta(s,n)^{-1}(x) = A_x \{0,1\}^\N.
        \end{equation}
        Since $h_\beta(s,n,\xf) \leq n$ for all $\xf \in \{0,1\}^\N$, we get that $A_x \subseteq \{0,1\}^{\leq n}$.
        \item Let $u, v \in A_x$, and assume that $\len u \leq \len v$. We will show by contradiction that $u \sqsubset v$. By definition of $A_x$, there exists $\uf,\vf \in \AM_\beta(s,n)^{-1}(x)$ such that $\uf_{1:h_\beta(s,n,\uf)} = u$ and $\vf_{1:h_\beta(s,n,\vf)} = v$. By sake of contradiction, assume that $u \not\sqsubset v$, i.e. there exists $i \in \{1,\ldots, \len u\}$, such that $u_{1:i-1} = v_{1:i-1}$, and $u_i \neq v_i$. Denote by $n_i$ the $i$-th element of $\SM_\beta(s,\uf)$ in increasing order. Then,
        \begin{equation}
            h_\beta(s,n_i,\uf) = \#\SM_\beta(s,n_i,\uf) = \# \left(\SM_\beta(s,\uf) \cap \{0, \ldots, n_i-1\} \right) = i-1,
        \end{equation}
        since by definition there are exactly $i$ elements in $\SM_\beta(s,\uf)$ that are not larger than $n_i$. Further note that
        \begin{equation}
            \uf_{1:h_\beta(s,n_i,\uf)} = \uf_{1:i-1} = u_{1:i-1} = v_{1:i-1} = \vf_{1:i-1} = \vf_{1:h_\beta(s,n_i,\uf)}.
        \end{equation}
        Therefore, $\vf$ is such that there exists $\yf \in \{0,1\}^\N$ satisfying $\vf = \uf_{1:h_\beta(s,n_i,\uf)} \yf$, so we can apply Lemma \ref{lem:the random beta expansion algorithm reads only the first bits of its input}. In particular,
        \begin{align}
            i-1 \overset{(a)}{=} h_\beta(s,n_i,\uf) &= h_\beta(s,n_i,\vf)\\
            r_i := \pi_1 K_\beta^{n_i}(s,\uf) &= \pi_1 K_\beta^{n_i}(s,\vf) \in S_\beta\\
            \pi_2 K_\beta^{n_i}(s,\uf) &= \sigma^{i-1}(\uf)\\
            \pi_2 K_\beta^{n_i}(s,\vf) &= \sigma^{i-1}(\vf).
        \end{align}
        where (a) follows by $n_i \in \SM_\beta(s,n_i,\uf)$. This further implies that 
        \begin{equation}
            \left[ \AM_\beta(s,n)(\uf)\right]_{n_i} = \tilde T_\beta\left( K_\beta^{n_i}(s,\uf)\right) = \tilde T_\beta \left( r_i, \sigma^{i-1}(\uf) \right) = \left[\sigma^{i-1}(\uf) \right]_i = \uf_i,
        \end{equation}
        and 
        \begin{equation}
            \left[ \AM_\beta(s,n)(\vf)\right]_{n_i} = \tilde T_\beta\left( K_\beta^{n_i}(s,\vf)\right) = \tilde T_\beta \left( r_i, \sigma^{i-1}(\vf) \right) = \left[\sigma^{i-1}(\vf) \right]_i = \vf_i.
        \end{equation}
        Since $\uf_i \neq \vf_i$, we get $\AM_\beta(s,n)(\uf) \neq \AM_\beta(s,n)(\vf)$, which is a contraction, since $\uf, \vf \in \AM_\beta(s,n)^{-1}(x)$. Therefore, $u \sqsubset v$.
        \item Since $A_x$ is a finite set, there exists an element of minimal length $u^\ast \in A_x$. Then, for every $v \in A_x$, $\len u^\ast \leq \len v$, so by the previous result, $u^\ast \sqsubset v$, which further implies $v \{0,1\}^\N \subseteq u^\ast \{0,1\}^\N$. This yields
        \begin{equation}
            \AM_\beta(s,n)^{-1}(x) = A_x \{0,1\}^\N = \bigcup_{v \in A_x} v \{0,1\}^\N = u^\ast \{0,1\}^\N.
        \end{equation}
        We define $w_\beta(s,x) := u^\ast$. Note that by definition of $A_x$, there exists $\xf^\ast \in \AM_\beta(s,n)^{-1}(x)$ such that $\xf^\ast_{1:h_\beta(s,n,\xf^\ast)} = w_\beta(s,x)$. Moreover, for every $\xf \in \AM_\beta(s,n)^{-1}(x)$, $\xf_{1:h_\beta(s,n,\xf)} \in A_x$. By the preceding point ,we get that $\xf^\ast_{1:h_\beta(s,n,\xf^\ast)} = w_\beta(s,x) \sqsubset \xf_{1:h_\beta(s,n,\xf)}$, i.e., there exists $\yf \in \{0,1\}^\N$ such that $\xf = \xf^\ast_{1:h_\beta(s,n,\xf^\ast)} \yf$. By Lemma \ref{lem:the random beta expansion algorithm reads only the first bits of its input},
        \begin{equation}
            h_\beta(s,n,\xf) = h_\beta(s,n,\xf^\ast) = |w_\beta(s,x)|.
        \end{equation}
        Finally, since $\AM_\beta(s,n)^{-1}(x) = w_\beta(s,x) \{0,1\}^\N$, then 
        \begin{equation}
            \xf_{1:h_\beta(s,n,\xf)} = \xf_{1:|w_\beta(s,x)|} = w_\beta(s,x).
        \end{equation}
        This concludes the proof.
    \end{enumerate}
\end{proof}

\noindent We now show how to infer this sequence of coin tosses from $x$ and $s$. We define
\begin{equation}
    \tilde \SM_\beta(s,x) = \bigcap_{\xf \in \AM_\beta(s,\len x)^{-1}(x)} \SM_\beta(s,\len x,\xf),
\end{equation}
and 
\begin{equation}
    \tilde h_\beta(s,x) := \# \tilde \SM_\beta(s,x).
\end{equation}
\begin{lemma}\label{lem:every preimage of a beta expansion throught the dajani algorithm yields the occurences in the switch set}
    Let $\beta \in (1,2)$, $s \in I_\beta$, and $x \in \{0,1\}^\ast$. Then, for all $\xf,\yf \in \AM_\beta(s,n)^{-1}(x)$,
    \begin{equation}
        \SM_\beta(s,\len x,\xf) = \SM_\beta(s,\len x,\yf).
    \end{equation}
    Consequently,
    \begin{equation}
        \tilde \SM_\beta(s,x) = \SM_\beta(s,\len x,\xf),
    \end{equation}
    for all $\xf \in  \AM_\beta(s,\len x)^{-1}(x)$.
\end{lemma}
\begin{proof}
    Let $\beta \in (1,2)$, $s \in I_\beta$, $x \in \{0,1\}^\ast$ and $\xf,\yf \in \AM_\beta(s,n)^{-1}(x)$. Then, by Lemma \ref{lem:the set of preimages by the dajani algorithm is a cylinder}, $\xf_{1:h_\beta(s,|x|,\xf)} = \yf_{1:h_\beta(s,|x|,\yf)} = w_\beta(s,x)$. Hence, there exists $\zf \in \{0,1\}^\N$ such that $\yf = w_\beta(s,x) \zf = \xf_{1:h_\beta(s,|x|,\xf)} \zf$. Then, by Lemma \ref{lem:the random beta expansion algorithm reads only the first bits of its input},
    \begin{equation}
        \SM_\beta(s,\len x,\xf) = \SM_\beta(s,\len x,\yf).
    \end{equation}
\end{proof}
% REFAIRE CETTE PREUVE AVEC LE LEMME
% \begin{proof}
%     Let $\beta \in (1,2)$, $s \in I_\beta$, and $x \in \{0,1\}^\ast$. Define $n := \len x$. Let $\xf, \yf \in \AM_\beta(s,n)^{-1}(x)$. By (\ref{eq:1:proof:lem:the set of preimages by the dajani algorithm is a cylinder}), we have 
%         \begin{equation}
%             \xf_{1:h_\beta(s,n,\xf)} \{0,1\}^\N \cup \yf_{1:h_\beta(s,n,\yf)} \{0,1\}^\N \subseteq \AM_\beta(s,n)^{-1}(x) = w_\beta(s,x) \{0,1\}^\N,
%         \end{equation}
%         hence either $\xf_{1:h_\beta(s,n,\xf)} \sqsubset \yf_{1:h_\beta(s,n,\yf)}$ or $\yf_{1:h_\beta(s,n,\yf)} \sqsubset \xf_{1:h_\beta(s,n,\xf)}$. W.l.o.g., we assume that $\xf_{1:h_\beta(s,n,\xf)} \sqsubset \yf_{1:h_\beta(s,n,\yf)}$, i.e., there exists $\zf \in \{0,1\}^\N$ such that $\yf = \xf_{1:h_\beta(s,n,\xf)} \zf$.
        
%         Then, $h_\beta(s,n,\yf) \geq h_\beta(s,n,\xf)$. Denote by $n_1 < \ldots < n_{h_\beta(s,n,\yf)}$ the elements of $\SM_\beta(s,n,\yf)$. We define 
%         \begin{equation}
%             N := \begin{cases}
%                 n_{h_\beta(s,n,\xf)} \ \text{if} \ h_\beta(s,n,\xf) = h_\beta(s,n,\yf),\\
%                 n_{h_\beta(s,n,\xf)} + 1 \ \text{if} \ h_\beta(s,n,\xf) < h_\beta(s,n,\yf).
%             \end{cases}
%         \end{equation}
%         We show that 
%         \begin{equation}\label{eq:0:proof:lem:every preimage of a beta expansion throught the dajani algorithm yields the occurences in the switch set}
%             \pi_1 K_\beta^i(s,\xf) = \pi_1 K_\beta^i(s,\yf), \ \forall i \leq N.
%         \end{equation}
%         First, $\pi_1 K_\beta^0(s,\xf) = \pi_1 K_\beta^0(s,\yf) = s$. Second, assume that for some $i \in \{1, \ldots, N-1\}$, $\pi_1 K_\beta^j(s,\xf) = \pi_1 K_\beta^j(s,\yf) =: r_j$ for all $j \leq i$. We show that $\pi_1 K_\beta^{i+1}(s,\xf) = \pi_1 K_\beta^{i+1}(s,\yf)$. 
%         \begin{enumerate}[label = (\alph*)]
%             \item First, note that
%             \begin{align}
%                 \SM_\beta(s,i-1,\xf)& = \{ j \in \{1,\ldots, i-1\} : \pi_1 K_\beta^j(s,\xf) \in S_\beta\}\\
%                 &= \{ j \in \{1,\ldots, i-1\} : \pi_1 K_\beta^j(s,\yf) \in S_\beta\}\\
%                 &= \SM_\beta(s,i-1,\yf) =: h_i.
%             \end{align}
%             Therefore,
%             \begin{equation}\label{eq:1:proof:lem:every preimage of a beta expansion throught the dajani algorithm yields the occurences in the switch set}
%                 \pi_2 K_\beta^i(s,\xf) = \xf_{h_i+1:\infty},
%             \end{equation}
%             and 
%             \begin{equation}\label{eq:2:proof:lem:every preimage of a beta expansion throught the dajani algorithm yields the occurences in the switch set}
%                 \pi_2 K_\beta^i(s,\xf) = \yf_{h_i+1:\infty}.
%             \end{equation}
%             \item \begin{enumerate}[label = (\roman*)]
%                 \item Suppose that $r_i \in E_\beta^{(0)} \cup E_\beta^{(1)}$, and define $b \in \{0,1\}$ such that $r_i \in E_\beta^{(b)}$. Then, 
%                 \begin{equation}
%                     \pi_1 K_\beta^{i+1}(s,\xf) = \pi_1 K_\beta(K_\beta^i(s,\xf)) = \beta r_i - b,
%                 \end{equation}
%                 and 
%                 \begin{equation}
%                     \pi_1 K_\beta^{i+1}(s,\yf) = \pi_1 K_\beta(K_\beta^i(s,\yf)) = \beta r_i - b,
%                 \end{equation}
%                 so 
%                 \begin{equation}
%                     \pi_1 K_\beta^{i+1}(s,\xf) = \pi_1 K_\beta^{i+1}(s,\yf).
%                 \end{equation}
%                 \item Suppose that $r_i \in S_\beta$. Then, by (\ref{eq:1:proof:lem:every preimage of a beta expansion throught the dajani algorithm yields the occurences in the switch set}) and (\ref{eq:2:proof:lem:every preimage of a beta expansion throught the dajani algorithm yields the occurences in the switch set}), we have
%                 \begin{equation}
%                     \pi_1 K_\beta^{i+1}(s,\xf) = \pi_1 K_\beta(r_i, \xf_{h_i+1:\infty}) = \beta r_i - \xf_{h_i+1},
%                 \end{equation}
%                 and 
%                 \begin{equation}
%                     \pi_1 K_\beta^{i+1}(s,\yf) = \pi_1 K_\beta(r_i, \yf_{h_i+1:\infty}) = \beta r_i - \yf_{h_i+1}.
%                 \end{equation}
%                 Since $\xf_{1:h_\beta(s,n,\xf)} \sqsubset \yf_{1:h_\beta(s,n,\yf)}$, we get that $\xf_{h_i+1} = \yf_{h_i+1}$, which yields
%                 \begin{equation}
%                     \pi_1 K_\beta^{i+1}(s,\xf) = \pi_1 K_\beta^{i+1}(s,\yf).
%                 \end{equation}
                
%             \end{enumerate}
%         \end{enumerate}
%         We then have proved (\ref{eq:0:proof:lem:every preimage of a beta expansion throught the dajani algorithm yields the occurences in the switch set}) by induction. We now end the proof.
%         \begin{enumerate}[label = (\alph*)]
%             \item We show that $h_\beta(s,n,\xf) = h_\beta(s,n,\yf)$. By sake of contradiction, assume that $h_\beta(s,n,\yf) > h_\beta(s,n,\xf)$. From (\ref{eq:0:proof:lem:every preimage of a beta expansion throught the dajani algorithm yields the occurences in the switch set}), we get that 
%             \begin{equation}
%                 \pi_1 K_\beta^{n_{h_\beta(s,n,\xf)+1}}(s,\xf) = \pi_1 K_\beta^{n_{h_\beta(s,n,\xf)+1}}(s,\yf) \overset{(a)}{\in} S_\beta,
%             \end{equation}
%             where (a) follows from $n_{h_\beta(s,n,\xf)+1} \in \SM_\beta(s,n,\yf)$. Therefore, $n_{h_\beta(s,n,\xf)+1} \in \SM_\beta(s,n,\xf)$, so $\# \SM_\beta(s,n,\xf) \geq h_\beta(s,n,\xf)+1$, which is a contradiction. 
%             \item By (\ref{eq:0:proof:lem:every preimage of a beta expansion throught the dajani algorithm yields the occurences in the switch set}), we have $\pi_1 K_\beta^i(s,\xf) = \pi_1 K_\beta^i(s,\yf)$ for all $i \in n_{h_\beta(s,n,\xf)} = n_{h_\beta(s,n,\yf)}$. 
%             Note that since $n_{h_\beta(s,n,\yf)} \leq n$ is the largest element of $\SM_\beta(s,n,\yf)$, we have
%             \begin{align}
%                 \SM_\beta(s,n,\yf) &= \SM_\beta(s,n,\yf) \cap \{1, \ldots, n_{h_\beta(s,n,\yf)}\}\\
%                 &=\{ i \in \{1,\ldots, n_{h_\beta(s,n,\yf)}\} : \pi_1 K_\beta^i(s,\yf) \in S_\beta \}\\
%                 &= \SM_\beta(s,n_{h_\beta(s,n,\yf)},\yf).
%             \end{align}
%             Further, we get that 
%             \begin{align}
%                 \SM_\beta(s,n,\yf) &= \SM_\beta(s,n_{h_\beta(s,n,\yf)},\yf)\\
%                 &= \{ i \in \{1,\ldots, n_{h_\beta(s,n,\yf)}\} : \pi_1 K_\beta^i(s,\yf) \in S_\beta \}\\
%                 &= \{ i \in \{1,\ldots, n_{h_\beta(s,n,\yf)}\} : \pi_1 K_\beta^i(s,\xf) \in S_\beta \}\\
%                 &= \SM_\beta(s,n_{h_\beta(s,n,\yf)},\xf)\\
%                 &\subseteq \SM_\beta(s,n,\xf).
%             \end{align}
%             So $\SM_\beta(s,n,\yf) \subseteq \SM_\beta(s,n,\xf)$, and moreover, as $h_\beta(s,n,\xf) = h_\beta(s,n,\yf)$, those two sets have the same number of elements, so $\SM_\beta(s,n,\yf) = \SM_\beta(s,n,\xf)$.
%         \end{enumerate}
% \end{proof}
The set $\tilde \SM_\beta(s,x)$ can be exploited to infer $w_\beta(s,x)$.
\begin{lemma}\label{lem:the noise can be expressed from the time steps for which the trajectory is in switch set}
    Let $\beta \in (1,2)$, $s \in I_\beta$ and $x \in \{0,1\}^\ast$. 
    \begin{equation}
        w_\beta(s,x) = \coprod \left\{ x_{i+1} : i \in \tilde \SM_\beta(s,x)\right\}.
    \end{equation}
\end{lemma}

\begin{proof}
    Let $\beta \in (1,2)$, $s \in I_\beta$, $x \in \{0,1\}^\ast$ and $\xf \in \AM_\beta(s,n)^{-1}(x)$. Then, 
    \begin{align}
        \left\{ x_{i+1} : i \in \tilde \SM_\beta(s,x)\right\} &= \left\{ x_{i+1} : i \in  \SM_\beta(s,|x|,\xf)\right\}\\
        &= \left\{\left[\AM_\beta(s,n)(\xf)\right]_{i+1} : i \in  \SM_\beta(s,|x|,\xf) \right\}\\
        &\overset{(\ref{eq:definition of the random beta expansion algorithm as a random variable})}{=} \left\{\tilde T_\beta\left(K_\beta^i(s,\xf) \right) : i \in  \SM_\beta(s,|x|,\xf) \right\}\\
        &=\left\{\tilde T_\beta\left(\pi_1K_\beta^i(s,\xf), \pi_2K_\beta^i(s,\xf) \right) : i \in  \SM_\beta(s,|x|,\xf) \right\}\\
        &\overset{(a)}{=} \left\{\left[\pi_2K_\beta^i(s,\xf)\right]_1 : i \in  \SM_\beta(s,|x|,\xf) \right\}\\
        &\osetlem{\ref{lem:the random beta expansion algorithm reads only the first bits of its input}}{=} \left\{\left[\sigma^{h_\beta(s,i,\xf)}(\xf)\right]_1 : i \in  \SM_\beta(s,|x|,\xf) \right\}\\
        &= \left\{\xf_{h_\beta(s,i,\xf)+1} : i \in  \SM_\beta(s,|x|,\xf) \right\}\\
        &\overset{(\ref{eq:main1:lem:recursive relation for determining h beta})}{=} \left\{\xf_{h_\beta(s,i+1,\xf)} : i \in  \SM_\beta(s,|x|,\xf) \right\}\\
        &\overset{(\ref{eq:main2:lem:recursive relation for determining h beta})}{=} \left\{\xf_j : j \in \{1,\ldots, h_\beta(s,|x|,\xf)\} \right\},
    \end{align}
    where (a) follows from $i \in \SM_\beta(s,|x|,\xf) \Rightarrow \pi_1K_\beta^i(s,\xf) \in S_\beta$, together with (\ref{eq:definition of the generalized random bit generator}).Hence,
    \begin{equation}
        \coprod \left\{ x_{i+1} : i \in \tilde \SM_\beta(s,|x|,\xf)\right\} = \coprod \left\{\xf_j : j \in \{1,\ldots, h_\beta(s,|x|,\xf)\} \right\} = \xf_{1:h_\beta(s,|x|,\xf)}.
    \end{equation}
    By Lemma \ref{lem:the set of preimages by the dajani algorithm is a cylinder}, $\xf_{1:h_\beta(s,|x|,\xf)} = w_\beta(s,x)$, which concludes the proof.
    
    % Then, by Lemma \ref{lem:every preimage of a beta expansion throught the dajani algorithm yields the occurences in the switch set}, $\tilde \SM_\beta(s,x) = \SM_\beta(s,\len x,\xf)$. We proceed by induction. If $\len x = 0$, then $x = \epsilon$, and $w_\beta(s,x) = \epsilon = \coprod \varnothing$. Now, suppose that EQUATION is satisfied for every $x \in \{0,1\}^n$, for some $n \in \N$. We prove that this is true for every $x \in \{0,1\}^{n+1}$. Let $x \in \{0,1\}^{n+1}$. Then,
    % \begin{equation}
    %     w_\beta(s,x_{1:n}) = \coprod \left\{ x_{i+1} : i \in \tilde \SM_\beta(s,x_{1:n})\right\}.
    % \end{equation}
    
    % Moreover, for $i \in \SM_\beta(s,\len x, \xf)$, $K_\beta^i(s,\xf) \in S_\beta$, so
    % \begin{equation}
    %     x_{i+1} = \tilde T_\beta\left(K_\beta^i(s,\xf) \right) = \left(\pi_2 K_\beta^i(s,\xf)\right)_1.
    % \end{equation}

\end{proof}

We now show how to derive $\tilde \SM_\beta(s,x)$ from $\Sigma_\beta(s,\len x)$ and $x$.
\begin{lemma}\label{lem:characterization of the set of indices in the switch set by divergences in the tree of beta expansions}
    Let $\beta \in (1,2)$, $s \in I_\beta$ and $x \in \{0,1\}^\ast$. Then,
    \begin{equation}
        \tilde \SM_\beta(s,x) = \left\{ i \in \{1,\ldots, \len x\}: \exists y \in \Sigma_\beta(s,\len x) \ s.t. \ y_{1:i-1} = x_{1:i-1}, y_i \neq x_i \right\}.
    \end{equation}
\end{lemma}
\begin{proof}
    Let $\beta \in (1,2)$, $s \in I_\beta$ and $x \in \{0,1\}^\ast$. Let 
    \begin{equation}
        A(s,x) := \left\{ i \in \{0,\ldots, \len x -1\}: \exists y \in \Sigma_\beta(s,\len x) \ s.t. \ y_{1:i} = x_{1:i}, y_{i+1} \neq x_{i+1} \right\}.
    \end{equation}
    We prove that $\tilde \SM_\beta(s,x) = A(s,x)$ by double inclusion.
    \begin{enumerate}[label = (\alph*)]
        \item We show that $\tilde \SM_\beta(s,x) \subseteq  A(s,x)$. Let $j \in \tilde \SM_\beta(s,x)$ and $\xf \in \AM_\beta(s,\len x)^{-1}(x)$. By Lemma \ref{lem:every preimage of a beta expansion throught the dajani algorithm yields the occurences in the switch set}, $\SM_\beta(s,\len x, \xf) = \tilde \SM_\beta(s,x)$, so $j \in \SM_\beta(s,\len x, \xf)$. Let $\yf \in \{0,1\}^\N$, such that $\yf_{1:h_\beta(s,j,\xf)} = \xf_{1:h_\beta(s,j,\xf)}$ and $\yf_{1:h_\beta(s,j,\xf)+1} \neq \xf_{1:h_\beta(s,j,\xf)+1}$, and define $y := \AM_\beta(s,\len x)(\yf)$. Note that $y \in \Sigma_\beta(s,|x|)$. We will prove that such a $y$ satisfies the condition in the definition of $A(s,x)$. By Lemma \ref{lem:the random beta expansion algorithm reads only the first bits of its input}, 
        \begin{align}
            \SM_\beta(s,i, \xf) &= \SM_\beta(s,i, \yf) =: \SM_i\\
            h_\beta(s,i,\xf) &= h_\beta(s,i,\yf) =: h_i\\
            \pi_1 K_\beta^i(s,\xf) &= \pi_1 K_\beta^i(s,\yf) =: r_i\\
            \pi_2 K_\beta^i(s,\xf) &= \sigma^{h_\beta(s,i,\xf)}(\xf) = \sigma^{h_i}(\xf)\\
            \pi_2 K_\beta^i(s,\yf) &= \sigma^{h_\beta(s,i,\yf)}(\yf) = \sigma^{h_i}(\yf),
        \end{align}
        for all $i \leq j$. Then,we have 
        \begin{align}
            x_{i+1} = \left[\AM_\beta(s,\len x)(\xf)\right]_{i+1} &= \tilde T_\beta \left(K_\beta^i(s,\xf) \right)\\
            &= \tilde T_\beta \left(r_i, \sigma^{h_i}(\xf)\right)\\
            &= \begin{cases}
                b \ \text{if} \ r_i \in E_\beta^{(b)}, \ b \in \{0,1\},\\
                \left[\sigma^{h_i}(\xf)\right]_1 \ \text{if} \ r_i \in S_\beta
            \end{cases}\\
            &= \label{eq:0:proof:lem:characterization of the set of indices in the switch set by divergences in the tree of beta expansions} \begin{cases}
                b \ \text{if} \ r_i \in E_\beta^{(b)}, \ b \in \{0,1\},\\
                \xf_{h_i+1} \ \text{if} \ r_i \in S_\beta,
            \end{cases}
        \end{align}
        for all $i \leq j$. By similar calculations, we obtain 
        \begin{equation}
            y_{i+1} = \begin{cases}
                b \ \text{if} \ r_i \in E_\beta^{(b)}, \ b \in \{0,1\},\\
                \yf_{h_i+1} \ \text{if} \ r_i \in S_\beta,
            \end{cases}
        \end{equation}
        for all $i \leq j$. Then, for $i \in \{0, \ldots, j-1\}$ such that $i \notin \SM_j$, we get 
        \begin{equation}\label{eq:1:proof:lem:characterization of the set of indices in the switch set by divergences in the tree of beta expansions}
            x_{i+1} = y_{i+1}.
        \end{equation}
        Moreover, for $i \in \{0, \ldots, j\}$ such that $i \in \SM_j$, we have
        \begin{equation}\label{eq:2:proof:lem:characterization of the set of indices in the switch set by divergences in the tree of beta expansions}
            x_{i+1} = \xf_{h_i+1} = \xf_{h_{i+1}}
        \end{equation}
        and 
        \begin{equation}\label{eq:3:proof:lem:characterization of the set of indices in the switch set by divergences in the tree of beta expansions}
            y_{i+1} = \yf_{h_i+1} = \yf_{h_{i+1}}.
        \end{equation}
        Since by assumption, $\xf_{h_{i+1}} = \yf_{h_{i+1}}$ for all $i < j$, and $\xf_{h_{j}+1} \neq \yf_{h_{j}+1}$, we get by combining (\ref{eq:1:proof:lem:characterization of the set of indices in the switch set by divergences in the tree of beta expansions}), (\ref{eq:2:proof:lem:characterization of the set of indices in the switch set by divergences in the tree of beta expansions}) and (\ref{eq:3:proof:lem:characterization of the set of indices in the switch set by divergences in the tree of beta expansions}) that
        \begin{equation}
            x_{1:j} = y_{1:j}, \ \text{and} \ x_{j+1} \neq y_{j+1}.
        \end{equation}
        Then, $x \in A(s,x)$. This concludes the proof that $\tilde \SM_\beta(s,x) \subseteq  A(s,x)$.
        \item We show that $\tilde \SM_\beta(s,x) \supseteq  A(s,x)$. Let $j \in A(s,x)$. Then, there exists $y \in \Sigma_\beta(s,\len x)$ such that 
        \begin{equation}
            x_{1:j} = y_{1:j}, \ \text{and} \ x_{j+1} \neq y_{j+1}.
        \end{equation}
        We prove that $j \in \tilde \SM_\beta(s,x)$ which is equivalent to prove that $j \in \SM_\beta(s,\len x, \xf)$ for some $\xf \in \{0,1\}^\N$. Let $\xf \in \AM_\beta(s,\len x)^{-1}(x)$ and $\yf \in \AM_\beta(s,\len x)^{-1}(y)$. In particular, $\xf \in \AM_\beta(s,j)^{-1}(x_{1:j})$ and $\yf \in \AM_\beta(s,j)^{-1}(y_{1:j}) = \AM_\beta(s,j)^{-1}(x_{1:j})$, so $w_\beta(s,x_{1:j})$ is a prefix of both $\xf$ and $\yf$, and by Lemma \ref{lem:the set of preimages by the dajani algorithm is a cylinder}, we have that 
        \begin{equation}
            \xf_{1:h_\beta(s,j,\xf)} = \yf_{1:h_\beta(s,j,\yf)} = w_\beta(s,x).
        \end{equation}
        By Lemma \ref{lem:the random beta expansion algorithm reads only the first bits of its input}, we have that
        \begin{align}
            h_\beta(s,j,\xf) &= h_\beta(s,j,\yf) =: h_j\\
            \pi_1 K_\beta^j(s,\xf) &= \pi_1 K_\beta^j(s,\yf) =: r_j\\
            \pi_2 K_\beta^j(s,\xf) &= \sigma^{h_\beta(s,j,\xf)}(\xf) = \sigma^{h_j}(\xf)\\
            \pi_2 K_\beta^j(s,\yf) &= \sigma^{h_\beta(s,j,\yf)}(\yf) = \sigma^{h_j}(\yf),
        \end{align}
        Then, following the calculations that lead to (\ref{eq:0:proof:lem:characterization of the set of indices in the switch set by divergences in the tree of beta expansions}),
        \begin{align}
            x_{j+1} = \begin{cases}
                b \ \text{if} \ r_j \in E_\beta^{(b)}, \ b \in \{0,1\},\\
                \xf_{h_j+1} \ \text{if} \ r_j \in S_\beta,
            \end{cases}
        \end{align}
        and 
        \begin{align}
            y_{j+1} = \begin{cases}
                b \ \text{if} \ r_j \in E_\beta^{(b)}, \ b \in \{0,1\},\\
                \yf_{h_j+1} \ \text{if} \ r_j \in S_\beta,
            \end{cases}
        \end{align}
    \end{enumerate}
    Suppose that $r_j \notin S_\beta$. Then, $x_{j+1} = y_{j+1}$, which is a contradiction. Then, $r_j \in S_\beta$, so $j \in \SM_\beta(s,\len x, \xf)$. This concludes the proof.
\end{proof}

We finally define a computable function that allows to calculate $w_\beta(s,x)$, given $x$ and the set of all $\beta$-expansions of $s$.
\begin{definition}
    Let $f_{1,all \to tosses} : \{0,1\}^\ast \to \{0,1\}^\ast$ to be defined such that 
    \begin{equation}
        f_{1,all \to tosses}\left(\langle x, \langle \Sigma_\beta(s,\len x) \rangle \rangle\right) = w_\beta(s,x),
    \end{equation}
    for every $\beta \in (1,2)$, $s \in I_\beta$ and $x \in \{0,1\}^\ast$.
\end{definition}
\begin{lemma}\label{lem:f 1 all to noise is computable}
    $f_{1,all \to tosses}$ is computable.
\end{lemma}
\begin{proof}
    We construct an algorithm that computes $f_{1,all \to tosses}$ as follows.
    \begin{algorithm}[ht]
        \caption{Algorithm for computing $f_{1,all \to tosses}$}
        \label{alg:algorithm that computes f 1 all noise}
        \begin{algorithmic}[1]
            \Require $x \in \{0,1\}^\ast$, $\langle \Sigma_\beta(s, \len x) \rangle$
            \State $n \gets \len x$. \label{line:1:alg:algorithm that computes f 1 all noise}
            \State Use $\langle \Sigma_\beta(s, n) \rangle$ to compute $k = \NN_\beta(s,n)$.\label{line:2:alg:algorithm that computes f 1 all noise}
            \State Use $\langle \Sigma_\beta(s, n) \rangle$ to generate all the elements $y \in \Sigma_\beta(s, n)$.\label{line:3:alg:algorithm that computes f 1 all noise}
            \State $S \gets \{\Sigma_\beta(s, n)\}$\label{line:4:alg:algorithm that computes f 1 all noise}
            \For{$i = 1, \ldots ,n$}\label{line:5:alg:algorithm that computes f 1 all noise}
                \If {there is $y \in \Sigma_\beta(s,n)$ such that $y_{1:i} = x_{1:i}$ and $y_{i+1} \neq x_{i+1}$}\label{line:7:alg:algorithm that computes f 1 all noise}
                    \State $S \gets S \cup \{i\}$.\label{line:8:alg:algorithm that computes f 1 all noise}
                \EndIf
            \EndFor
            \State \Return $\coprod S$.\label{line:12:alg:algorithm that computes f 1 all noise}
        \end{algorithmic}
    \end{algorithm}

    \noindent Note that the set $S$ constructed by the algorithm exactly satisfies 
    \begin{equation}
        S = \left\{ i \in \{1,\ldots, \len x\}: \exists y \in \Sigma_\beta(s,\len x) \ s.t. \ y_{1:i-1} = x_{1:i-1}, y_i \neq x_i \right\}.
    \end{equation}
    By combining Lemma \ref{lem:the noise can be expressed from the time steps for which the trajectory is in switch set} and Lemma \ref{lem:characterization of the set of indices in the switch set by divergences in the tree of beta expansions}, we get that indeed
    \begin{equation}
        w_\beta(s,x) = \coprod S.
    \end{equation}
\end{proof}

\subsection{\texorpdfstring{Distribution of algorithmic complexity of $\beta$-expansions}{}}
We are now ready to move on to the proofs on algorithmic complexity. We combine the different functions defined previously to build a new multivalued function.
\begin{definition}\label{def:f beta to 2 + noise}
    Let $\beta \in (1,2)$, $f_{\beta \to tosses} : \{0,1\}^\ast \rightrightarrows \{0,1\}^\ast$ be defined as 
    \begin{equation}
        f_{\beta \to tosses}(x) = \left\{ f_{\beta,1+all\to tosses}(\langle x, y\rangle) : y \in f_{\beta, 1 \to all}(x)\right\},
    \end{equation}
    and $f_{\beta \to 2+tosses} : \{0,1\}^\ast \rightrightarrows \{0,1\}^\ast$ be defined by 
    \begin{equation}\label{eq:main:def:f beta to 2 + noise}
        f_{\beta \to 2+tosses}(x) = \left\{ \left\langle y,z \right\rangle : (y,z) \in f_{\beta \to 2}(x) \times f_{\beta \to tosses}(x) \right\},
    \end{equation}
    for all $x \in \{0,1\}^\ast$.
\end{definition}
\begin{lemma}
    Let $\beta \in (1,2)$, $s \in [0,1]$, $n \in \N$ and $\xf \in \Sigma_\beta(s)$. Then,
    \begin{equation}
        \langle \yf_{1:n}, w_\beta(s,\xf_{1:\nlogbeta})\rangle \in f_{\beta \to 2+tosses}(\xf_{1:\nlogbeta}),
    \end{equation}
    for all $n \in \N$, where $\yf$ is the greedy binary expansion of $s$.
\end{lemma}
\begin{proof}
    Let $\beta \in (1,2)$, $s \in [0,1]$, $n \in \N$, $\xf \in \Sigma_\beta(s)$ and define $\yf$ to be the greedy binary expansion of $s$. Then,
    \begin{equation}
        \yf_{1:n} \in f_{\beta \to 2}(\xf_{1:\nlogbeta}),
    \end{equation}
    and 
    \begin{equation}
        \langle \Sigma_\beta(s, \nlogbeta) \rangle \in f_{\beta, 1 \to all}(\xf_{1:\nlogbeta}),
    \end{equation}
    so 
    \begin{equation}
        w_\beta(s, \xf_{1:\nlogbeta}) = f_{\beta,1+all\to tosses}(\langle x, \langle \Sigma_\beta(s,\nlogbeta)\rangle\rangle) \in f_{\beta \to tosses}(\xf_{1:\nlogbeta}).
    \end{equation}
    Then,
    \begin{equation}
        (\yf_{1:n}, w_\beta(s, \xf_{1:\nlogbeta})) \in f_{\beta \to 2}(\xf_{1:\nlogbeta}) \times f_{\beta \to tosses}(\xf_{1:\nlogbeta}),
    \end{equation}
    and the Lemma follows from (\ref{eq:main:def:f beta to 2 + noise}).
\end{proof}
\begin{lemma}\label{lem:size of the multifunction converting a beta expansion into its noise and the binary expansion}
    Let $\beta \in (1,2)$ be an algebraic number. Then, for all $n \in \N$, and all $x \in \{0,1\}^n$,
    \begin{equation}
        \# f_{\beta \to 2+tosses}(x) \leq \left(\frac{1}{\beta-1} + 3\right) \left(\frac{1}{\beta-1} + 2\right)^2\frac{n^{2k_\beta}}{(L_\beta\Pi_\beta)^2}\left(L_\beta \Pi_\beta^+ \right)^{2n}.
    \end{equation}
\end{lemma}
\begin{proof}
    Let $\beta \in (1,2)$ be an algebraic number, and $x \in \{0,1\}^\ast$. First, by Lemma \ref{lem:upper bound on the cardinality of the beta to 2 conversion multifunction},
    \begin{equation}
        \#f_{\beta \to 2}(x) \leq \frac{1}{\beta-1} + 3.
    \end{equation}
    Second, by Lemma \ref{lem:size of the multifunction that generates all beta expansions from one},
    \begin{equation}
        \#f_{\beta, 1 \to all}(x) \leq \left(\frac{1}{\beta-1} + 2\right)^2\frac{n^{2k_\beta}}{(L_\beta\Pi_\beta)^2}\left(L_\beta \Pi_\beta^+ \right)^{2n},
    \end{equation}
    and therefore 
    \begin{equation}
        \#f_{\beta \to tosses}(x) \leq \left(\frac{1}{\beta-1} + 2\right)^2\frac{n^{2k_\beta}}{(L_\beta\Pi_\beta)^2}\left(L_\beta \Pi_\beta^+ \right)^{2n}.
    \end{equation}
    The proof is concluded by noting that
    \begin{align}
        \# f_{\beta \to 2+tosses}(x) &= \#f_{\beta \to 2}(x)  \times \#f_{\beta \to tosses}(x).
    \end{align}
\end{proof}
We now establish the subsequent results on algorithmic complexity. The equations are very involved, and we first make the following remark. Algebraic numbers are computable, i.e., for every algebraic number $\beta$, there exists a computable function $\varphi: \N \to \Q$ such that $\limi{n} \varphi(n) = \beta$. A consequence of this is that 
\begin{equation}
    K[x|\beta] = K[x] + \bigOx[1]{|x|}, \ \forall x \in \{0,1\}^\ast.
\end{equation}
In what follows, we focus on $\beta$ being algebraic, and we hence replace $K[\cdot |\beta]$ by $K[\cdot]$.
\begin{theorem}\label{thm:lower bound of kolmogorov complexity for beta expansions when beta is algebraic}
    Let $\beta \in (1,2)$ be an algebraic number, $s \in [0,1]$ and $n \in \N$. Denote by $\yf \in \{0,1\}^\N$ the greedy binary expansion of $s$. Then,
    \begin{align}\label{eq:algebraic:thm:lower bound of kolmogorov complexity for beta expansions when beta is algebraic}
        K[\langle \yf_{1:n}, w_\beta(s,\xf_{1:\nlogbeta})\rangle] &\leq K[\xf_{1:\nlogbeta}] + 2n \log_\beta\left(L_\beta \Pi_\beta^+\right)
        + \bigOx[\log n]{n},
    \end{align}
for all $\xf \in \Sigma_\beta(s)$. Moreover, if $\beta$ is a Pisot number,
\begin{equation}\label{eq:pisot:thm:lower bound of kolmogorov complexity for beta expansions when beta is algebraic}
    K[\langle \yf_{1:n}, w_\beta(s,\xf_{1:\nlogbeta})\rangle] \leq K[\xf_{1:\nlogbeta}] + \bigOx[1]{n},
\end{equation}
for all $\xf \in \Sigma_\beta(s)$. 
\end{theorem}
\begin{proof}
    Let $\beta \in (1,2)$ be an algebraic number, $s \in [0,1]$, $n \in \N$ and $\xf \in \Sigma_\beta(s)$. Then, by combination of 
    \begin{equation}
        \langle \yf_{1:n}, w_\beta(s,\xf_{1:\nlogbeta})\rangle \in f_{\beta \to 2+tosses}(\xf_{1:\nlogbeta}).
    \end{equation}
    Moreover, $f_{\beta \to 2+tosses}$ is computable. Hence, by Theorem \ref{thm:relatively computable multivalued functions and complexity},
    \begin{align}
        K[\langle \yf_{1:n}, w_\beta(s,\xf_{1:\nlogbeta})\rangle] &\leq K[\xf_{1:\nlogbeta}] + \log \#f_{\beta \to 2+tosses}(\xf_{1:\nlogbeta})\\
        & + \log \log \#f_{\beta \to 2+tosses}(\xf_{1:\nlogbeta}) + \bigOx[1]{n}.
    \end{align}
    Note that by Lemma \ref{lem:size of the multifunction converting a beta expansion into its noise and the binary expansion}, 
    \begin{equation}
        \log \#f_{\beta \to 2+tosses}(\xf_{1:\nlogbeta}) \leq 2n \log_\beta \left(L_\beta \Pi^+_\beta \right) + 2k_\beta \log_\beta n + \bigOx[1]{n}
    \end{equation}
    and 
    \begin{equation}
        \log \log \#f_{\beta \to 2+tosses}(\xf_{1:\nlogbeta}) \leq \log n + \log \log_\beta n + \bigOx[1]{n},
    \end{equation}
    hence (\ref{eq:algebraic:thm:lower bound of kolmogorov complexity for beta expansions when beta is algebraic}) follows. If specifically, $\beta$ is a Pisot number, then $L_\beta = \Pi^+_\beta = 1$ and $k_\beta = 0$, so 
    \begin{equation}
        \log \#f_{\beta \to 2+tosses}(\xf_{1:\nlogbeta}) \leq \bigOx[1]{n}
    \end{equation}
    and hence
    \begin{equation}
        \log \log \#f_{\beta \to 2+tosses}(\xf_{1:\nlogbeta})\leq \bigOx[1]{n}.
    \end{equation}
    (\ref{eq:pisot:thm:lower bound of kolmogorov complexity for beta expansions when beta is algebraic}) follows.
\end{proof}
\noindent In what follows, we use  
\begin{equation}
    K[x|y] := K[x|y0^\infty], \ \forall x,y \in \{0,1\}^\ast.
\end{equation}
\begin{corollary}\label{cor:lower bound of kolmogorov complexity for beta expansions when beta is algebraic}
    Let $\beta \in (1,2)$ be an algebraic number, $s \in [0,1]$ and $n \in \N$. Denote by $\yf \in \{0,1\}^\N$ the greedy binary expansion of $s$. Then,
    \begin{align}
        K[\xf_{1:\nlogbeta}] &\geq K[\yf_{1:n}] + K[w_\beta(s,\xf_{1:\nlogbeta})|\yf_{1:n}] - 2n \log_\beta\left(L_\beta \Pi_\beta^+\right) \nonumber\\ 
        &+ \bigOx[\log n]{n},\label{eq:main:cor:lower bound of kolmogorov complexity for beta expansions when beta is algebraic}
    \end{align}
    for all $\xf \in \Sigma_\beta(s)$.
\end{corollary}
\begin{proof}
    Let $\beta \in (1,2)$ be an algebraic number, $s \in [0,1]$ and $n \in \N$.
    By \cite[Theorem 2.8.2]{li2008introduction}, 
    \begin{equation}
        K[\langle y,x\rangle] \geq K[y] + K[x|y] + \bigOw[\log |\langle y,x\rangle|]{|\langle y,x \rangle|},
    \end{equation}
    for all $x,y \in \{0,1\}^\ast$. Therefore, in particular,
    \begin{equation}\label{eq:0:proof:cor:lower bound of kolmogorov complexity for beta expansions when beta is algebraic}
        K[\langle y,x\rangle] \geq K[y] + K[x|y] + \bigOx[\log n]{n},
    \end{equation}
    for all $x \in \{0,1\}^n$ and $y \in \{0,1\}^{\leq \nlogbeta}$. (\ref{eq:main:cor:lower bound of kolmogorov complexity for beta expansions when beta is algebraic}) follows immediately by combining (\ref{eq:0:proof:cor:lower bound of kolmogorov complexity for beta expansions when beta is algebraic}) and Theorem \ref{thm:lower bound of kolmogorov complexity for beta expansions when beta is algebraic}.
\end{proof}
\noindent We define a distribution of the Kolmogorov complexity of $\beta$-expansions. Let 
\begin{equation}
    \KK_\beta[s,n,k] := \left\{ x \in \Sigma_\beta(s,\nlogbeta) : K[x] \leq K[\yf_{1:n}] + k - 2n \log_\beta\left(L_\beta \Pi_\beta^+\right)\right\}.
\end{equation}

\begin{corollary}
    Let $\beta \in (1,2)$ be an algebraic number, $s \in [0,1]$ and $n,k \in \N$. Then, there exists $M \in \N$ such that
    \begin{equation}
        \# \KK_\beta[s,n,k] \leq 2^{k+1} n^M.
    \end{equation}
\end{corollary}
\begin{proof}
    Let $\beta \in (1,2)$ be an algebraic number, $s \in [0,1]$ and $n,k \in \N$. By Corollary \ref{cor:lower bound of kolmogorov complexity for beta expansions when beta is algebraic}, there exists $M \in \N$ such that 
    \begin{align}
        K[\xf_{1:\nlogbeta}] &\geq K[\yf_{1:n}] + K[w_\beta(s,\xf_{1:\nlogbeta})|\yf_{1:n}] - 2n \log_\beta\left(L_\beta \Pi_\beta^+\right)\\ 
        &- M\log n,
    \end{align}
    for all $\xf \in \Sigma_\beta(s)$. Define the set 
    \begin{equation}
        \WW_\beta[s,n,k] := \left\{ x \in \Sigma_\beta(s,\nlogbeta) : K[w_\beta(s,x)|\yf_{1:n}] \leq  k - M\log n\right\}.
    \end{equation}
    Then, by Corollary \ref{cor:lower bound of kolmogorov complexity for beta expansions when beta is algebraic}, we have 
    \begin{equation}
        \KK_\beta[s,n,k] \subseteq \WW_\beta[s,n,k].
    \end{equation}
    For $z \in \{0,1\}^\ast$, let $z^\ast$ be the canonical sequence for $z$ with respect to $\yf_{1:n}$, i.e., the lexicographically maximal sequence satisfying $|z^\ast| = K[z|\yf_{1:n}]$. Note that $z \mapsto z^\ast$ injective. Note that, by injectivity of $x\mapsto w_\beta(s,x)$, we have
    \begin{align}
        \#\WW_\beta[s,n,k] &= \#\left\{ x \in \Sigma_\beta(s,\nlogbeta) : K[w_\beta(s,x)|\yf_{1:n}] \leq  k + M\log n\right\}\\
        &\overset{(a)}{=} \#\left\{ w_\beta(s,x) : x \in \Sigma_\beta(s,\nlogbeta),\right.\\
        & \hspace{0.45in} \left. |K[w_\beta(s,x)|\yf_{1:n}] \leq  k + M\log n\right\}\\
        &\leq \#\left\{ z \in \{0,1\}^\ast : K[z|\yf_{1:n}] \leq k + M \log n\right\}\\
        &\overset{(b)}{\leq} \{ z^\ast \in  : z \in \{0,1\}^\ast,  |z^\ast| \leq k + M \log n\}\\
        &= \# \{0,1\}^{\leq k + M \log n} = \sum_{i=0}^{k + M\log n} 2^i\\
        & = 2^{k + M \log n+1} - 1 \leq 2^{k+1} n^M,
    \end{align}
    where (a) is by injectivity of $x \mapsto w_\beta(s,x)$, (b) follows by injectivity of $z \mapsto z^\ast$.
\end{proof}